\section{上下文工程：帮助 Agent 理解环境}

\subsection{Context Engineering 的核心理念}

Harness Engineering 的一个核心职能是做好\textbf{上下文的交付工程}（Delivery Mechanism for Context Engineering）。Agent 在未知环境中执行任务时，面临大量上下文缺失：目录结构是什么？有哪些工具可用？评估标准是什么？时间限制多少？

\begin{importantbox}{Context Engineering 核心原则}
Agent 对其环境、约束和评估标准了解得越充分，就越能自主地高质量完成工作。Harness 工程师的职责本质是：\textbf{代替 Agent 准备和交付上下文，使 Agent 能自主完成工作}。
\end{importantbox}

\subsection{三类上下文注入}

LangChain 实施了三类关键的上下文注入：

\subsubsection{目录结构与工具发现}

LocalContextMiddleware 在 Agent 启动时自动执行：
\begin{itemize}[leftmargin=1.5em]
    \item 映射当前工作目录及其父目录、子目录的结构
    \item 运行 Bash 命令发现可用工具（如 Python 安装路径）
    \item 将发现结果注入 Agent 上下文
\end{itemize}

这一步至关重要，因为上下文发现和搜索\textbf{本身就容易出错}。让 Agent 自己去探索环境会引入不必要的错误风险。通过确定性注入，可以消除这一错误面。

\subsubsection{可测试性提示}

Agent 通常不知道自己的代码需要被程序化测试验证。LangChain 在 Prompt 中明确告知：
\begin{itemize}[leftmargin=1.5em]
    \item 你的代码将被自动化测试脚本评估（类似于提交代码后的 CI）
    \item Task Spec 中提到的文件路径必须严格遵循
    \item 不仅要处理 Happy Path，还要考虑 Edge Cases
\end{itemize}

\begin{warningbox}{"Slop Buildup"现象}
如果不提示 Agent 写可测试的代码，Agent 倾向于只验证"一切正常"的路径（Happy Path），忽略边界条件。随着时间推移，这种松散的编码习惯会导致"Slop Buildup"——代码在表面上看起来能工作，但在严格测试下会大量失败。强制模型遵循测试标准是一个有效的反制策略。
\end{warningbox}

\subsubsection{时间预算管理}

Agent 在时间估计方面表现很差。LangChain 注入时间预算警告（Time Budget Warnings），在剩余时间不足时提醒 Agent：
\begin{itemize}[leftmargin=1.5em]
    \item 停止探索新方向
    \item 完成当前工作
    \item 转入验证阶段
\end{itemize}

虽然真实世界的编程通常没有严格时限，但在基准测试（和很多生产场景）中，如果不告知 Agent 时间约束，它会一直"探索"而无法收敛。

\subsection{本章小结}

Context Engineering 的核心思想是"不要让 Agent 去猜"。通过在 Harness 层面主动注入目录结构、工具信息、可测试性要求和时间约束，可以显著减少 Agent 因信息不足而犯错的概率。这是一种"为 Agent 做 Onboarding"的策略。


